\section{Постановка задачи}
\label{sec:task}

Целью лабораторной работы является изучение механизма ответа на сообщения в системе NIKA: проектирование логических правил и фраз, в том числе с шаблонами, зависящими от состояния базы знаний.

Вариант~\mdash  15. Предметная область~\mdash  музыка. Классы сообщений \texttt{about\_performer}, \texttt{about\_author} и \texttt{about\_album} заданы в лабораторной работе №2.

Требуется:
\begin{enumerate}
	\item спроектировать логические правила на каждый класс фраз ответа;
	\item в посылках правил использовать и класс сообщения, и выделенную сущность;
	\item показать разные ответы на одно и то же намерение в зависимости от состояния БЗ, используя шаблонные фразы;
	\item получить ответ в пользовательском интерфейсе NIKA;
	\item описать фразы: без повторов, не менее двух фраз в каждом классе, не менее трёх шаблонных фраз на весь диалог.
\end{enumerate}
